\section{Build notes}
\label{sec:build}

\subsection{Configure \& build (Windows)}
The repository-root \pathref{build.py} program is a standard-library front end
over the existing CMake targets. It prints the commands it runs and does not
replace or delete CMake configuration. With Python on \texttt{PATH}, open its
interactive terminal build console with:
\begin{lstlisting}[language={}]
python build.py
\end{lstlisting}

The arrow keys select the configuration, optional documentation and Tracy,
parallelism, and an action. The same operations remain explicit commands for
scripts and CI. For example, a Debug build is:
\begin{lstlisting}[language={}]
python build.py build --config Debug
\end{lstlisting}
Focused examples are:
\begin{lstlisting}[language={}]
python build.py build --config Debug --target Illumo --parallel
python build.py test
python build.py test --test Illumo.CellGame.SaveLoadRoundTrip
python build.py run -- -ww 1280 -wh 720
\end{lstlisting}
With redirected standard input or output, omitting the command performs the
normal Release all-target build instead of opening the console.
Use \texttt{--no-docs} to opt out of the optional PDF target,
\texttt{--build-dir} to select another tree, and \texttt{--dry-run} to inspect
the commands. The program deliberately has no clean command.

The dashboard's \textbf{Run existing build} action, or
\texttt{run --no-build}, skips configuration and compilation and fails if the
selected executable does not exist:
\begin{lstlisting}[language={}]
python build.py run --config Debug --no-build
\end{lstlisting}
The normal \texttt{run} command still builds before launching. Direct CMake
remains supported:
\begin{lstlisting}[language={}]
cmake -S Illumo -B build -G "Visual Studio 17 2022" -A x64
cmake --build build --config Release
\end{lstlisting}

The default build compiles the application and \texttt{IllumoTests}, then runs
every granular \texttt{Illumo} CTest entry through the
\texttt{IllumoRunTests} target. Use an explicit target only for a focused
build.

\symbolref{DebugModule} is intentionally absent from Release. CMake adds
\pathref{Source/Engine/DebugModule.cpp} only for the Debug configuration, and
\symbolref{CellMain} includes/registers the module only when
\texttt{NDEBUG} is not defined (D-B1).

\begin{lstlisting}[language={}]
cmake --build build --config Release --target IllumoTests
\end{lstlisting}

Artifacts:
\begin{itemize}
  \item \pathref{build/Release/Illumo.exe}
  \item \pathref{build/Debug/Illumo.exe}
  \item \pathref{build/Release/IllumoTests.exe} (convenience CMake target
        \texttt{IllumoRenderTests} aliases \texttt{IllumoTests})
\end{itemize}

\subsection{Tests (headless)}
The Python front end configures and builds \texttt{IllumoTests} before invoking
CTest, or invokes one exact runner case when \texttt{--test} is supplied:
\begin{lstlisting}[language={}]
python build.py test
python build.py test --list-tests
python build.py test --test Illumo.CellGame.SaveLoadRoundTrip
\end{lstlisting}
The underlying commands remain available:
\begin{lstlisting}[language={}]
ctest --test-dir build -C Release -N -L Illumo
ctest --test-dir build -C Release -L Illumo --output-on-failure
build/Release/IllumoTests.exe --list
build/Release/IllumoTests.exe --run Illumo.CellGame.SaveLoadRoundTrip
\end{lstlisting}
One runner avoids recompiling shared production sources, but each logical case
is a separately named CTest process with its own working directory (D-T1). No
OpenGL context is required (\texttt{MockBackend}).

\clearpage
\subsection{Coverage (Clang/LLVM)}
Run the complete coverage workflow with:
\begin{lstlisting}[language={}]
python build.py coverage
\end{lstlisting}
This composes the following existing configure and target commands:
\begin{lstlisting}[language={}]
cmake -S Illumo -B build-coverage -G Ninja \
  -DCMAKE_BUILD_TYPE=Debug \
  -DCMAKE_C_COMPILER=clang -DCMAKE_CXX_COMPILER=clang++ \
  -DILLUMO_BUILD_DOCUMENTATION=OFF -DILLUMO_ENABLE_COVERAGE=ON
cmake --build build-coverage --target IllumoCoverage
\end{lstlisting}
The target runs every granular test, merges per-process profiles, prints an
LLVM report, fails below 85\% production line coverage, and writes
\pathref{build-coverage/coverage-html/index.html}. The denominator excludes
tests, vendored/system code, MockBackend, and the live OpenGL backend. Platform
dialogs/windows and live OpenGL remain manual smoke-test responsibilities.

\subsection{Documentation PDF}
From the repository root, run:
\begin{lstlisting}[language={}]
python build.py docs
\end{lstlisting}
This delegates to the existing PowerShell script, which can also be called
directly:
\begin{lstlisting}[language={}]
.\docs\build.ps1
\end{lstlisting}
The script creates \pathref{docs/output/} and builds two artifacts:
\begin{itemize}
  \item \texttt{illumo.pdf} from \pathref{docs/latex/illumo.tex};
  \item \texttt{architecture-map.pdf} from
        \pathref{docs/latex/architecture-map.tex}.
\end{itemize}
On Windows, a normal all-target CMake build invokes the same script through
\texttt{IllumoDocs} when PowerShell and \texttt{latexmk} were found at configure
time. To opt out, configure with:
\begin{lstlisting}[language={}]
-DILLUMO_BUILD_DOCUMENTATION=OFF
\end{lstlisting}

\subsection{Known setup issues}
\begin{itemize}
  \item GLEW is built directly from
        \pathref{Illumo/thirdparty/glew-2.1.0/src/glew.c}; the upstream
        optional CMake subtree is not required.
  \item Debug may enable \texttt{/fsanitize=address}; requires a toolchain that
        supports MSVC ASan.
  \item FreeType may warn about missing optional system deps (HarfBuzz, zlib,
        PNG); build can still succeed with bundled defaults.
  \item GLSL under \pathref{Illumo/Shader/} is copied beside each application
        and test executable by a post-build command.
\end{itemize}

\subsection{House style (CONTRIBUTING)}
See \pathref{docs/contributing.md}: limited \texttt{auto}, avoid namespaces,
no recursion, third-party via PR.
